% =====================================================================
%  PFE — Chapter: Agent Orchestration & Post-Audit Optimization
%
%  Drop-in chapter for the author's PFE book. All figures are produced
%  by docs/pfe/03-figures/scripts/build_all.py and saved as PDF/PGF in
%  docs/pfe/03-figures/. All cited measurements come from
%  docs/pfe/02-optimizations/<n>.md.
%
%  Required preamble in the host book (excerpt):
%    \usepackage{graphicx}
%    \usepackage{pgfplots}
%    \usepackage{booktabs}
%    \usepackage{listings}
%    \usepackage{hyperref}
% =====================================================================

\chapter{Agent Orchestration and Post-Audit Optimization}
\label{ch:opt}

% ---------------------------------------------------------------------
\section{Introduction}
% ---------------------------------------------------------------------

This chapter recounts the engineering work done on FINAL\_MPFE, an AI-orchestrated
syllabus and worksheet generation platform, after its first full release. The
first 69 pull requests of the project (Phase 0 through Phase 6, summarised in
Appendix~\ref{app:history}) brought the system from an empty repository to a
three-agent platform with grounded generation, anchored UX, and a typed
streaming protocol. A deep audit, performed against a running instance, then
identified four classes of integration bug that this chapter resolves.

The chapter is organised around the audit's priority order. Each subsequent
section presents one optimisation: the problem statement, the root-cause
analysis, the design alternatives that were considered, the chosen design, the
measurement methodology, and the before/after numbers.

% ---------------------------------------------------------------------
\section{System under study}
% ---------------------------------------------------------------------

% Pull a one-paragraph architecture summary here from 00-history.md
% Phase 0 — Phase 6 invariants. Reference Figure 1: agent topology.

\begin{figure}[h]
  \centering
  % \includegraphics[width=0.85\linewidth]{03-figures/agent-topology.pdf}
  \caption{Agent topology of FINAL\_MPFE post-PR\#91. Three agents
  (syllabus generator, activity-tooled, activity-toolless) share a
  supervisor router and emit typed state slices over the Vercel AI SDK
  v5 UI Message Stream protocol (SSE with typed JSON frames). The
  activity-tooled agent grounds its output via an MCP server backed by
  Supabase.}
  \label{fig:topology}
\end{figure}

% ---------------------------------------------------------------------
\section{The audit}
\label{sec:audit}
% ---------------------------------------------------------------------

% Summarise 01-audit.md TL;DR in 4 paragraphs:
%   1. Method (read + run + measure)
%   2. The four classes of integration bug
%   3. Priority ordering rationale
%   4. Forward pointer to optimisations
% Reference Table 1: audit findings by severity.

% ---------------------------------------------------------------------
\section{Optimisation 1 — State hydration on completion}
\label{sec:opt-hydration}
% ---------------------------------------------------------------------
% Source: docs/pfe/02-optimizations/01-hydration.md
% Figures: hydration-before-after.pdf

% ---------------------------------------------------------------------
\section{Optimisation 2 — Parallel research with LangGraph \texttt{Send}}
\label{sec:opt-parallel}
% ---------------------------------------------------------------------
% Source: docs/pfe/02-optimizations/02-parallel-research.md
% Figures: research-latency-cdf.pdf, research-fanout-diagram.pdf
%
% Key claims for the chapter:
%   - Sequentialism in async-friendly code is a recurring anti-pattern in
%     LangGraph applications because the framework's per-step reducer
%     model nudges developers towards state-machine traversal.
%   - The Send API restores the fan-out semantics without sacrificing
%     checkpointer durability.
%   - Real measurements: P50/P95/P99 latency before vs after, on the same
%     5-topic test set, with and without scrape failures injected.

% ---------------------------------------------------------------------
\section{Optimisation 3 — Tier rebalancing for the critic}
\label{sec:opt-critic}
% ---------------------------------------------------------------------
% Source: docs/pfe/02-optimizations/03-critic-tier.md
% Figures: tier-token-cost.pdf, critic-tier-quality-comparison.pdf

% ---------------------------------------------------------------------
\section{Optimisation 4 — Intake message hygiene}
\label{sec:opt-intake}
% ---------------------------------------------------------------------
% Source: docs/pfe/02-optimizations/04-intake-dedup.md
% No quantitative figure; UX before/after screenshots in the PR.

% ---------------------------------------------------------------------
\section{Optimisation 5 — Source-aware research card}
\label{sec:opt-sources}
% ---------------------------------------------------------------------
% Source: docs/pfe/02-optimizations/05-research-sources.md

% ---------------------------------------------------------------------
\section{Optimisation 6 — Polish bundle}
\label{sec:opt-polish}
% ---------------------------------------------------------------------
% Source: docs/pfe/02-optimizations/06-quick-wins.md

% ---------------------------------------------------------------------
\section{Optimisation 7 — Writer recovery}
\label{sec:opt-writer-recovery}
% ---------------------------------------------------------------------
% Source: docs/pfe/02-optimizations/07-writer-recovery.md

% ---------------------------------------------------------------------
\section{Optimisation 8 — Frontend perf: ETag polling and memo deps}
\label{sec:opt-fe-perf}
% ---------------------------------------------------------------------
% Source: docs/pfe/02-optimizations/08-fe-perf.md

% ---------------------------------------------------------------------
\section{Optimisation 9 — Live LLM token streaming}
\label{sec:opt-token-streaming}
% ---------------------------------------------------------------------
% Source: docs/pfe/02-optimizations/09-token-streaming.md
%
% Key claims for the chapter:
%   - Token-level streaming on the supervisor required `streamEvents`,
%     a partial-JSON byte walker, and an envelope reorder so the
%     user-visible `user_message` field is emitted first.
%   - On the FE, dropping `useDeferredValue` was necessary because
%     React's concurrent scheduler coalesces rapid source updates into
%     a single deferred render at end-of-stream.
%   - Time-to-first-token: 3.2 s P50 → 0.41 s P50 on the same prompt.

% ---------------------------------------------------------------------
\section{Optimisation 10 — AI SDK v4 → v5 wire migration}
\label{sec:opt-aisdk-v5}
% ---------------------------------------------------------------------
% Source: docs/pfe/02-optimizations/10-ai-sdk-v5.md
%
% Two production-grade ordering bugs (BUG-0001 — terminal run slice
% dropped by the v5 writer's internal `finished` flag; BUG-0002 —
% missing [DONE] terminator on the POST error path) caught by the
% post-merge review pass and fixed in the same PR. Both characteristic
% of v5 migrations — worth recording for the Discussion.

% ---------------------------------------------------------------------
\section{Discussion}
% ---------------------------------------------------------------------
% Three reflective threads:
%   1. The cost of integration vs the cost of components.
%   2. LangGraph's reducer model: a feature, not a constraint.
%   3. What measurement enabled and what it could not see.

% ---------------------------------------------------------------------
\appendix
\section{Pre-audit development history}
\label{app:history}
% ---------------------------------------------------------------------
% Pull from 00-history.md — phase summaries + the PR citation table.
